\begin{lemma}
If $r\geq 0$ and $\rho\geq 3/2$, then $G\left(\omega\right)$ in \eqref{eq:trancedental_eq} satisfies
\[
G\left(\omega\right)\leq c\rho=\pi,\qquad 0<\omega<\frac{c}{\pi m+r}.
\]
\end{lemma}

\begin{proof}
By Lemma \ref{lem:client-derivative-criterion}, the denominator of $G\left(\omega\right)$ is positive. Normalize the sampling rate by setting
\[
x=\frac{\pi m+r}{c}\omega,\qquad
a=\frac{cm}{\pi m+r}.
\]
Then $0<x<1$, $m\omega=ax$, and, since $r\geq 0$ and $\pi=c\rho$,
\[
0<a\leq \frac{1}{\rho}\leq \frac{2}{3}.
\]
In particular, $1-m\omega=1-ax>0$ throughout the interval.

The nonlinear term in the numerator of $G\left(\omega\right)$ can be written as
\[
\begin{aligned}
&(1-m\omega)^2\omega
\left(\frac{c}{(\pi m+r)\omega}\right)^{\frac{1}{1-m\omega}}\\
&\qquad=
\frac{c(1-ax)^2}{\pi m+r}\exp\left(
\frac{ax}{1-ax}\ln\frac{1}{x}
\right).
\end{aligned}
\]
Let $z=\frac{ax}{1-ax}\ln(1/x)$. Since $x\ln(1/x)\leq 1-x$,
\[
0\leq z\leq \frac{a(1-x)}{1-ax}\leq a\leq \frac{2}{3}.
\]
For $0\leq z\leq 2/3$, the exponential series gives $e^z\leq 1+2z$. Therefore,
\[
\begin{aligned}
&(1-m\omega)^2\omega
\left(\frac{c}{(\pi m+r)\omega}\right)^{\frac{1}{1-m\omega}}\\
&\qquad\leq
\frac{c(1-ax)^2}{\pi m+r}
\left(1+\frac{2ax}{1-ax}\ln\frac{1}{x}\right).
\end{aligned}
\]

Substituting this bound into the numerator of $G\left(\omega\right)$, and using the positivity of its denominator, shows that it suffices to prove
\[
\begin{aligned}
1-ax-\ln\frac{1}{x}
+\frac{c(1-ax)^2}{\pi m+r}
\left(1+\frac{2ax}{1-ax}\ln\frac{1}{x}\right)
\leq
\rho\left(1-ax-ax\ln\frac{1}{x}\right).
\end{aligned}
\]
Write $L=\ln(1/x)$. Since $c/(\pi m+r)=a/m$, the right-hand side minus the left-hand side is
\[
F_m=(\rho-1)(1-ax)+L(1-\rho ax)
-\frac{a}{m}(1-ax)^2-\frac{2a^2x}{m}(1-ax)L.
\]
As $m\geq 1$, we have $F_m\geq F$, where
\[
F=(\rho-1)(1-ax)+L(1-\rho ax)
-a(1-ax)^2-2a^2x(1-ax)L.
\]
Thus it remains to show that $F\geq 0$.

The coefficient of $\rho$ in $F$ is $1-ax(1+L)$. Because $x(1+\ln(1/x))\leq 1$ and $a\leq 2/3$, this coefficient is positive. Hence $F$ is bounded below by its value at $\rho=3/2$. At that value, write
\[
F=C(a)+B(a)L,
\]
where
\[
C(a)=\frac{1}{2}(1-ax)-a(1-ax)^2,\qquad
B(a)=1-\frac{3}{2}ax-2a^2x(1-ax).
\]
Both $C(a)$ and $B(a)$ decrease with $a$ on $(0,2/3]$. Indeed,
\[
C'(a)=-1-\frac{x}{2}+4ax-3a^2x^2
\leq -\frac{4}{3}\left(x-\frac{13}{16}\right)^2-\frac{23}{192}<0,
\]
and
\[
B'(a)=x\left(-\frac{3}{2}-4a+6a^2x\right)
\leq x\left(-\frac{3}{2}+2a(3a-2)\right)<0.
\]
Since $L>0$, replacing $a$ by $2/3$ therefore gives a lower bound:
\[
F\geq F\left(\frac{3}{2},\frac{2}{3},x\right)
=
\frac{(3-2x)(8x-3)}{54}
+(1-x)\ln\frac{1}{x}
-\frac{8}{9}\left(1-\frac{2x}{3}\right)x\ln\frac{1}{x}.
\]
For $0<x<1$, we have $\ln(1/x)\geq 1-x$ and
$x\ln(1/x)\leq (1-x^2)/2$. The latter follows by integrating
$1/t\leq(1+t^{-2})/2$ from $x$ to $1$. Applying these bounds yields
\[
\begin{aligned}
F
&\geq \frac{(3-2x)(8x-3)}{54}
+(1-x)^2
-\frac{4}{9}\left(1-\frac{2x}{3}\right)(1-x^2)\\
&=\frac{14\left(x-\frac{1}{2}\right)^2+16(1-x)^3+\frac{3}{2}}{54}>0.
\end{aligned}
\]
Thus $F_m\geq F>0$, which proves the required inequality and hence
$G\left(\omega\right)\leq c\rho=\pi$.
\end{proof}
